\documentclass[11pt]{article}
\usepackage[margin=1in]{geometry}
\usepackage{amsmath,amssymb,amsthm}
\usepackage[hidelinks]{hyperref}
\newtheorem{theorem}{Theorem}
\newcommand{\R}{\mathbb R}
\newcommand{\Obs}{\mathcal O}
\title{A complete negatively curved observation Gramian\protect\\with two-state ambiguity}
\author{Alec Kriebel\protect\\Independent researcher\protect\\
\small ORCID: \href{https://orcid.org/0009-0001-9320-500X}{0009-0001-9320-500X}}
\date{Version 1.0, 5 October 2026}
\begin{document}
\maketitle
\begin{abstract}
We give a smooth autonomous observed system on a connected closed genus-three
surface whose exact finite-time differential observation Gramian is complete
and has constant negative sectional curvature, although every attained output
history corresponds to exactly two initial states. The system has zero dynamics
and a Euclidean-valued output obtained by composing a connected double cover of
a closed hyperbolic surface with a smooth Nash isometric embedding into
$\R^{17}$. For each fixed $T>0$ its Gramian is $P_T=T\pi^*g$ and its curvature is
$-1/T$. Uniform positive upper and lower bounds hold in a fixed finite
normalized atlas and relative to any fixed smooth background metric. This
refutes the manifold implication printed in Krener's 2009 open problem under
those coherent meanings of its coordinate bounds. The construction is an application of classical constructions. A bounded
primary-source audit found no earlier explicit complete construction in the
inspected corpus; no absolute historical firstness is asserted.
\end{abstract}

\section{The printed problem and the claim}
Krener's contribution \cite{krener} considers $\dot x=f(x)$, $y=h(x)$,
with local coordinates on a state manifold and output in $\R^p$. For a fixed
$T>0$, global observability means injectivity of the initial-state history map
\[
 \Obs_T(x)(t)=h(\varphi_t(x)),\qquad 0\leq t<T,
\]
where $\varphi_t$ is the flow. Local observability in that contribution means
local injectivity of this map. Its exact differential Gramian is
\begin{equation}\label{eq:gramian}
 P_{T,x}(v,w)=\int_0^T
 \left\langle d(h\circ\varphi_t)_x v,
 d(h\circ\varphi_t)_x w\right\rangle\,dt.
\end{equation}
In local coordinates this is the printed integral
$\int_0^T\Phi(t)^\mathsf{T}H(t)^\mathsf{T}H(t)\Phi(t)\,dt$,
where $\Phi=d\varphi_t$ and $H=dh$ along the trajectory. The conjectured
implication is that uniform positive definiteness and boundedness of this
Gramian, uniform negative curvature, and geodesic completeness imply global
observability. The complete contribution, printed pp.~674--675, imposes no
simple-connectivity assumption, global $\R^n$ state chart, nonzero-dynamics
condition, or restriction on the finite output dimension.

The source leaves normalization for the matrix bounds unspecified. We verify
both common bounds in a fixed finite atlas and tensor bounds against any fixed
smooth background metric. Common coefficient bounds over arbitrary rescalings
of every possible chart cannot hold for any positive metric: under $z\mapsto az$
its matrix scales by $a^{-2}$. Such an additional quantifier would make the
hypothesis empty. Our claim concerns the printed manifold formulation with
fixed normalization. It does not settle a separately strengthened question
restricted to a single Euclidean state domain.

\begin{theorem}\label{thm:counterexample}
There exist a connected closed oriented genus-three smooth surface $M$, a
complete smooth vector field $f$ on $M$, and a smooth map $h:M\to\R^{17}$ such
that, for every fixed $T>0$:
\begin{enumerate}
\item $P_T$ in \eqref{eq:gramian} is a smooth positive definite, geodesically
complete metric of constant sectional curvature $-1/T$;
\item for any fixed smooth Riemannian metric $q$ on $M$ there are constants
$0<c\leq C<\infty$ with $Tcq\leq P_T\leq TCq$;
\item there is a fixed finite atlas in which
$4T I_2\leq [P_T]\leq (64T/9)I_2$;
\item $\Obs_T$ is locally injective, and every history in $\Obs_T(M)$ has exactly
two distinct initial states.
\end{enumerate}
The same pair of states is indistinguishable for all real times.
\end{theorem}

\section{Construction and proof}
\begin{proof}
Let $(S,g)$ be a closed oriented hyperbolic genus-two surface with curvature
$-1$. One concrete construction is a regular geodesic octagon of interior
angle $\pi/4$, with opposite sides identified by orientation-preserving
hyperbolic isometries reversing their boundary orientations. Number the corners
$v_i$ and sides $[v_i,v_{i+1}]$ modulo eight. The pairing identifies
$v_i$ with $v_{i+5}$ and $v_{i+1}$ with $v_{i+4}$ for $i=0,1,2,3$.
The vertex link is the single cycle
\[
 0\longrightarrow5\longrightarrow2\longrightarrow7\longrightarrow4
 \longrightarrow1\longrightarrow6\longrightarrow3\longrightarrow0.
\]
The eight sectors have total angle $2\pi$. Developing them successively
into the hyperbolic plane, the final transition fixes the vertex and has
derivative the full $2\pi$ rotation. An isometry fixing a point and its tangent
frame is the identity, so there is no residual vertex holonomy. The sectors
therefore form an ordinary hyperbolic disk. Across a paired edge two half disks
also form a smooth hyperbolic disk.
Thus the quotient metric is smooth at edges and at the vertex, without a cone
point. There is one vertex, four edges and one face, giving $\chi(S)=-2$ and
genus two. For completeness, the octagon exists in the Poincar\'e disk with
vertices $2^{-1/4}\exp(ik\pi/4)$: the right triangle from its center to a vertex
and side midpoint has angles $\pi/8,\pi/8,\pi/2$ and circumradius $R$ with
$\cosh R=\cot^2(\pi/8)=3+2\sqrt2$.

The surface-group presentation is
\[
 \pi_1(S)=\langle a_1,b_1,a_2,b_2\mid[a_1,b_1][a_2,b_2]=1\rangle.
\]
These are abstract standard surface generators, not asserted to be the
opposite-side polygon letters. Send $a_1$ to $1\in\mathbb Z/2$ and all other
generators to zero. The relation is respected because its commutators vanish
in this abelian group. The resulting homomorphism is surjective, with
index-two kernel. Covering-space classification \cite[Prop.~1.32,
Prop.~1.36 and Thm.~1.38]{hatcher} gives a connected two-sheeted cover
$\pi:M\to S$. Lifting the smooth charts gives a smooth local diffeomorphism
and the pulled-back orientation. This is an unramified cover, with two
distinct points in every fiber, rather than two disconnected copies of $S$.

The surface $M$ is compact. Indeed, finitely many compact sets lying inside
evenly covered neighborhoods can be chosen with interiors covering $S$;
the inverse image of each is the union of two compact copies. A lifted finite
cell structure gives $\chi(M)=2\chi(S)=-4$, so $M$ has genus three. Put
$\widetilde g=\pi^*g$. The map $\pi$ is a local isometry for this smooth metric,
and hence $K(\widetilde g)=-1$.

Nash's smooth compact isometric embedding theorem \cite[Thm.~2, p.~59]{nash}
supplies a smooth embedding
\[
 e:(S,g)\hookrightarrow\R^{17},\qquad e^*\langle\ ,\ \rangle=g.
\]
Its stated bound $n(3n+11)/2$ is $17$ when $n=2$, and the theorem includes
$C^\infty$ metrics and embeddings. We use the smooth embedding theorem,
including global injectivity, rather than only a $C^1$ immersion theorem.

Set $f\equiv0$ and $h=e\circ\pi$. The flow is $\varphi_t=\operatorname{id}_M$
for all $t\in\R$, so $d\varphi_t=\operatorname{id}$ and
\begin{align}
 P_{T,x}(v,w)
 &=T\langle dh_xv,dh_xw\rangle\notag\\
 &=Tg_{\pi(x)}(d\pi_xv,d\pi_xw)
 =T\widetilde g_x(v,w).\label{eq:identity}
\end{align}
In a fixed chart containing the constant trajectory, $Df=0$, $\Phi=I$ and
$H=Dh(x)$, giving exactly the same printed Gramian. Thus the negatively curved
metric is the actual observation metric, rather than an independently chosen
auxiliary metric. The endpoint convention $[0,T)$ has no effect on this integral.

Positive constant scaling leaves the Levi--Civita connection unchanged and
divides sectional curvature by the scale. Hence $K(P_T)=-1/T$. Each $P_T$ is
geodesically complete: on a compact Riemannian manifold a fixed-speed geodesic
stays in a compact tangent sphere bundle; smooth geodesic-flow extension then
prevents any finite-time endpoint. This also handles negative times; zero-speed
geodesics are constant. At $T=1$ the curvature is identically $-1$.

Fix a smooth background metric $q$. Its unit tangent bundle is compact, and
$\widetilde g_x(v,v)$ is continuous and strictly positive on it. Its attained
minimum $c>0$ and maximum $C<\infty$, followed by homogeneity, give
$Tcq\leq P_T\leq TCq$. For the explicit coordinate bounds, center sufficiently
small hyperbolic disk charts at the origin and restrict their coordinate radii
to at most $1/2$. Compactness supplies a finite subcover. In each such chart,
with $s=u^2+v^2\leq1/4$,
\[
 [P_T]=\frac{4T}{(1-s)^2}I_2,
 \qquad 4T I_2\leq[P_T]\leq\frac{64T}{9}I_2.
\]
This is an all-point inequality in each selected chart; it does not require a
radius-$1/2$ disk to inject at every point. The lower and upper differences at
$T=1$ are respectively
\[
 \frac{4s(2-s)}{(1-s)^2},\qquad
 \frac{4(1-4s)(7-4s)}{9(1-s)^2},
\]
both nonnegative throughout the stated interval.

On one sheet above a sufficiently small evenly covered neighborhood, $\pi$ is
injective, as is $e$. Thus $h$ and the constant-history map $\Obs_T$ are locally
injective. Conversely, for $T>0$,
\[
 \Obs_T(x)=\Obs_T(x')
 \quad\Longleftrightarrow\quad h(x)=h(x')
 \quad\Longleftrightarrow\quad\pi(x)=\pi(x').
\]
Every fiber of $\pi$ has exactly two distinct points. Therefore every
\emph{attained} history has exactly two initial states, while histories outside
$\Obs_T(M)$ have none. This proves all assertions.
\end{proof}

Taking $T=1$ gives a counterexample to the printed geometric implication.
Constants here are uniform over states for a fixed positive observation time.
There is no bound uniform over all $T\downarrow0$ or $T\uparrow\infty$:
$P_0=0$ and $K(P_T)\to0$ as $T\to\infty$.

\section{Context, priority and verification}
The proof is a classical synthesis of covering spaces, hyperbolic surface
geometry and Nash's embedding theorem. The broader distinction between
local differential tests and global state ambiguity is old. Hermann--Krener
\cite[Ex.~3.3, p.~734]{hk} give $\dot x=u$, $y=(\cos x,\sin x)$ with full
local rank and globally indistinguishable states $x+2k\pi$; their terminology
for local observability should not be conflated with the history-injectivity
definition above. Gauthier \cite[\S7.1, Thms.~9--10]{gauthier} treats covering
realizations in a complete controllable analytic framework. That framework
does not apply verbatim to our zero dynamics, and neither citation supplies
the complete nonvacuous two-dimensional negative-Gramian construction.
In dimension one, quantification over sectional two-planes can be empty; no
claim about the literal one-dimensional reading is needed here.

The original 2009 open-conjecture presentation has been authenticated. A
bounded primary-source audit, including complete readings of the ECC 2009
paper of Dirr--Helmke--Jordan \cite{dirr} and the 2013 survey coauthored by
Krener \cite{survey}, found no earlier explicit complete construction in the
inspected corpus. Neither complete source presents a resolution of the printed
curvature question. Dirr--Helmke--Jordan's Theorem~1.2 assumes that different
$\omega$-limit sets have different output-image sets (A3). For $f=0$ the limit
sets are singletons, and two different lifts have the same output; thus A3
fails. Their auxiliary complete metric is not the exact observation Gramian.
The survey treats the nonlinear empirical Gramian as local sensitivity and
uses additional injective-coordinate or system-structure hypotheses for
observer designs. Neither result entails the printed curvature-only implication.

Sussmann's 1974 quotient announcement \cite{sussmann} was also read in full;
its quotient-regularity and fibre-map criteria contain no observation-Gramian
curvature statement. A complete reading of Krishnaprasad's 1977 Allerton paper
\cite{allerton} concerns metrics on parameter/realization spaces, rather than
the state finite-history Gramian here. Later observation/bundle literature and
other relevant sources were compared at the scopes recorded in
\texttt{SOURCE\_AUDIT.md}. Some older realization originals and portions of the
worldwide literature remain uninspected. These are disclosed coverage limits,
and no identified source currently supplies affirmative evidence of the full
two-dimensional result. The negative search finding does not establish
worldwide continued openness, absolute firstness, or novelty of the covering,
Nash or local/global ambiguity mechanism. In particular we make no claim to
the first refutation of every literal dimensional reading of the conjecture.

The accompanying standard-library verifier checks elementary exact polynomial,
rational, monodromy and matrix identities, with a separate finite-precision
decimal diagnostic. It rejects optimized Python execution. Its finite checks
supplement the proof; they do not construct the Nash embedding or
machine-certify covering existence, smoothness or completeness. Those global
claims follow from the proof and the explicitly imported classical theorems.

\paragraph{AI use and review status.}
AI tools were used extensively in solving, drafting, reproducing and
adversarially reviewing this work. This is an unrefereed preprint without conventional human peer review.
Independent AI audits examined the proof and its source and priority scope;
they do not replace conventional human peer review.

\begin{thebibliography}{9}
\bibitem{krener} A.~J. Krener, Open Problem: Global Observability of Spaces of
Negative Curvature, Oberwolfach Report 11/2009, pp.~674--675.
\href{https://doi.org/10.4171/OWR/2009/11}{doi:10.4171/OWR/2009/11}.
\bibitem{nash} J. Nash, The Imbedding Problem for Riemannian Manifolds,
\emph{Annals of Mathematics} (2) \textbf{63} (1956), 20--63, Theorem~2, p.~59.
\href{https://sites.math.rutgers.edu/~feehan/teaching/math866/nash.pdf}{Original paper}.
\bibitem{hatcher} A. Hatcher, \emph{Algebraic Topology}, Cambridge University
Press, 2002, pp.~51, 61, 66--67.
\href{https://pi.math.cornell.edu/~hatcher/AT/AT.pdf}{Author's online version}.
\bibitem{hk} R. Hermann and A.~J. Krener, Nonlinear Controllability and
Observability, \emph{IEEE Transactions on Automatic Control} \textbf{22}
(1977), 728--740. \href{https://doi.org/10.1109/TAC.1977.1101601}{doi:10.1109/TAC.1977.1101601}.
\bibitem{gauthier} J.-P. Gauthier, Observability (deterministic systems) and
realization theory, author manuscript dated
29 June 2009, \S7.1. \href{https://pageperso.lis-lab.fr/jean-paul.gauthier/papers/115.pdf}{Author manuscript}.
Its original public posting date has not been verified.
\bibitem{dirr} G. Dirr, U. Helmke and J. Jordan, Global observability of real
analytic systems, \emph{European Control Conference} (2009), 4582--4586.
\href{https://doi.org/10.23919/ECC.2009.7075123}{doi:10.23919/ECC.2009.7075123}.

\bibitem{survey} W. Kang, A.~J. Krener, M. Xiao and L. Xu, A Survey of Observers
for Nonlinear Dynamical Systems, in \emph{Data Assimilation for Atmospheric,
Oceanic and Hydrologic Applications}, Vol.~II (2013), 1--25.
\href{https://doi.org/10.1007/978-3-642-35088-7_1}{doi:10.1007/978-3-642-35088-7\_1}.

\bibitem{sussmann} H.~J. Sussmann, On quotients of manifolds: a generalization
of the closed subgroup theorem, \emph{Bulletin of the American Mathematical
Society} \textbf{80} (1974), 573--575.
\href{https://doi.org/10.1090/S0002-9904-1974-13502-0}{doi:10.1090/S0002-9904-1974-13502-0}.
\bibitem{allerton} P.~S. Krishnaprasad, Geometry of Parametric Models: Some
Probabilistic Questions, \emph{Proceedings of the 15th Allerton Conference on
Control, Communication and Computing} (1977), 661--670.
\href{https://user.eng.umd.edu/~krishna/images/allerton_1977.pdf}{Author-hosted original}.
\end{thebibliography}
\end{document}
